\documentclass[a4paper,11pt]{article}
\usepackage[utf8]{inputenc}
\usepackage[T1]{fontenc}
\usepackage[french]{babel}
\usepackage{amsmath,amssymb}
\usepackage{enumitem}
\usepackage{booktabs}
\usepackage{array}
\usepackage{geometry}
\geometry{margin=2.2cm}

\title{Améliorations du modèle RL-PDNN \\ Intégration de contraintes réalistes pour l'IoT}
\author{}
\date{}

\begin{document}

\maketitle

\section{Contraintes proposées et justifications détaillées}

\subsection{Comparaison des deux approches d’intégration de l’énergie}

On compare ici en profondeur les deux façons classiques d’intégrer la contrainte énergétique dans un problème d’optimisation d’allocation de couches en IoT :

\begin{enumerate}[label=\arabic*.]
    \item \textbf{Énergie comme contrainte stricte} (hard constraint)
    \item \textbf{Énergie dans la fonction objectif} (approche multi-objectif scalarisée)
\end{enumerate}

\subsubsection{Option 1 – Énergie comme contrainte stricte}

\textbf{Formulation mathématique}
\[
\min \quad \text{Latence}
\]
sous la contrainte
\begin{equation}
\sum_{r^* \in RQ} \sum_{l=1}^{L_{r^*}} A^i_{r^*,l} \cdot E^l_{r^*} \le \bar{E}_i \quad \forall \, I_i
\label{eq:energie_contrainte}
\end{equation}

où :
\begin{itemize}[leftmargin=*]
    \item $E^l_{r^*}$ \quad = énergie nécessaire pour exécuter la couche $l$ du modèle $r^*$
    \item $\bar{E}_i$ \quad = budget énergétique disponible (maximum autorisé) du dispositif $I_i$
    \item $A^i_{r^*,l}$ \quad = variable de décision binaire indiquant si la couche $l$ du modèle $r^*$ est assignée au dispositif $i$
    \item $RQ$ \quad = ensemble des requêtes / modèles à inférer
    \item $L_{r^*}$ \quad = nombre total de couches du modèle $r^*$
\end{itemize}

\textbf{Signification}
\begin{itemize}[leftmargin=*]
    \item L’énergie est un \textbf{budget maximal} à ne jamais dépasser
    \item Toute solution qui respecte la limite est valide
    \item L’objectif principal reste \textbf{uniquement la latence}
\end{itemize}

\textbf{Exemple de comportement}
\begin{center}
\begin{tabular}{lcc}
\toprule
Solution & Latence (s) & Énergie (unités) \\
\midrule
A & 1.0 & 95 \\
B & 1.2 & 40 \\
\bottomrule
\end{tabular}
\end{center}

Budget = 100 unités $\Rightarrow$ les deux solutions sont valides.  
$\Rightarrow$ Le modèle choisit systématiquement \textbf{A} (plus rapide), sans se soucier de l’énergie restante.

\textbf{Impact théorique}
\begin{itemize}[leftmargin=*]
    \item Réduction de l’espace des solutions admissibles
    \item Risque de rendre le problème \textbf{infaisable} si le budget est trop restrictif
    \item Aucun compromis énergétique n’est recherché
\end{itemize}

\subsubsection{Option 2 – Énergie dans la fonction objectif (multi-objectif)}

\textbf{Formulation mathématique}
\[
\min \Bigl( \text{Latence} + \alpha \times \text{Énergie} \Bigr)
\]
(où $\alpha > 0$ est un coefficient de pondération)

\textbf{Signification}
L’énergie devient un \textbf{critère de performance} à part entière.  
On recherche activement un compromis entre :
\begin{itemize}[leftmargin=*]
    \item minimiser la latence (vitesse)
    \item minimiser la consommation énergétique
\end{itemize}

\textbf{Exemple de comportement ($\alpha = 0.5$)}
\begin{center}
\begin{tabular}{lccc}
\toprule
Solution & Latence (s) & Énergie & Valeur objectif \\
\midrule
A & 1.0 & 95 & 1.0 + 0.5×95 = 48.5 \\
B & 1.2 & 40 & 1.2 + 0.5×40 = 21.2 \\
\bottomrule
\end{tabular}
\end{center}

$\Rightarrow$ Le modèle choisit \textbf{B} (meilleure valeur globale).

\textbf{Impact théorique}
\begin{itemize}[leftmargin=*]
    \item Problème multi-objectif scalarisé
    \item Navigation possible sur le front de Pareto (via variation de $\alpha$)
    \item Le problème reste toujours faisable (pas de contrainte dure)
\end{itemize}

\subsubsection{Différence conceptuelle profonde}

\begin{itemize}[leftmargin=*]
    \item \textbf{Contrainte stricte} $\quad\leftrightarrow\quad$ \textbf{Logique physique / sécurité} \\
    \quad « La batterie ne peut physiquement pas dépasser cette limite. »
    
    \item \textbf{Fonction objectif} $\quad\leftrightarrow\quad$ \textbf{Logique économique / optimisation} \\
    \quad « Parmi les solutions possibles, je veux consommer le moins d’énergie possible. »
\end{itemize}

Le choix entre les deux dépend fortement du contexte :
\begin{itemize}[leftmargin=*]
    \item Applications critiques où dépasser le budget = panne certaine $\to$ contrainte stricte
    \item Applications longue durée où l’on veut maximiser l’autonomie $\to$ fonction objectif
\end{itemize}


\subsection{Contrainte de latence maximale (QoS)}
\textbf{Problème actuel}\\
Minimiser la latence $\neq$ garantir une latence acceptable.\\
Exemple : le RL trouve 1,5 s mais l’application exige $<$ 1 s $\rightarrow$ le modèle original ne garantit rien.

\textbf{Amélioration apportée}
\begin{enumerate}[label=\alph*)]
\item \textbf{Garantie de service (QoS)} : élimination des solutions invalides
\[
\text{Latency}_{r^*} \le L_{\max}
\]
Le système devient réellement utilisable en pratique.
\item \textbf{Applicabilité temps réel} : permet d’adresser
\begin{itemize}[leftmargin=*]
\item Vidéosurveillance
\item Santé connectée
\item Véhicules autonomes
\end{itemize}
$\Rightarrow$ très forte valeur industrielle.
\end{enumerate}

\subsection{Contrainte de confiance (Trust-aware allocation)}
\textbf{Problème actuel}\\
Le modèle protège contre l’exposition complète, mais considère implicitement que tous les devices sont équivalents. En réalité : certains sont compromis, d’autres fiables, d’autres publics.

\textbf{Amélioration apportée}
\begin{enumerate}[label=\alph*)]
\item \textbf{Sécurité adaptative} : les couches sensibles vont aux devices fiables
\[
A^i_{r^*,l} \le \text{Trust}_i
\]
\item \textbf{Défense contre attaques internes} : même dans le réseau, un device ne peut pas exécuter les couches critiques.
\item \textbf{Compatible Zero-Trust} : approche très moderne en cybersécurité $\Rightarrow$ grande valeur académique.
\end{enumerate}

\subsection{Contrainte de concurrence (capacité simultanée)}
\textbf{Problème actuel}\\
Le modèle ignore que les devices traitent plusieurs requêtes en parallèle.

\textbf{Amélioration apportée}
\begin{enumerate}[label=\alph*)]
\item \textbf{Gestion des files d’attente} : évite la surcharge simultanée.
\item \textbf{Réduction des délais cachés} : la latence réelle inclut le temps d’attente.
\end{enumerate}

\textbf{Formulation mathématique}
\[
\sum_{r^*} A^i_{r^*,l} \le K_i
\]

\subsection{Contrainte de robustesse (tolérance aux pannes)}
\textbf{Problème actuel}\\
Si un device tombe en panne $\rightarrow$ l’inférence échoue complètement.

\textbf{Amélioration apportée}
\begin{enumerate}[label=\alph*)]
\item \textbf{Fiabilité du système} : possibilité de re-routage.
\item \textbf{Haute disponibilité} : critique pour les systèmes santé, sécurité, etc.
\end{enumerate}

\end{document}